\section{Death and Stupidity}

Originally published in July, 2007 for a private seminar.

\begin{quotex}
Call no man happy until he is dead.
\flright{\textit{Ancient Greek proverb}}

Society cannot exist unless a controlling power upon will and appetite be placed somewhere, and the less of it there is within, the more of it there must be without.
\flright{\textsc{Edmund Burke}}
\end{quotex}


The question was recently asked: What factors contribute to keeping man in bondage to evil?

There are two traditional schools of thought. One school, represented by Socrates for example, emphasizes the Intellect: once a man is educated in the truth, his moral actions will then follow. Another school, Schopenhauer being a representative, emphasizes the Will: there is something inherently perverse about man's will that no amount of education can eradicate.

The Traditionalist position sees both as interrelated. The Christian teaching, for example, claims that man, such as he is, comes into the world with a weakened will and a darkened intellect. And although the ``truth will make you free'', the solution to life cannot be found in life, but must come from beyond life, or as a ``grace''. Let's explore how to understand that metaphysically.

There are two fundamental obstacles in life, whose absurdity and intractability permeate every aspect of life, making a mockery of every humanistic attempt at self-improvement or social development. These are \textbf{Death} and \textbf{Stupidity}.

Stupidity can take several forms such as ignorance of the facts, or holding onto beliefs that are illogical or metaphysically absurd. The solution, then, is obvious: study and observe the facts, learn the rules of logic, be guided by true metaphysical doctrines. Nothing in these steps is hidden and all are available to anyone diligent enough to seek them out.

Since stupidity is a defect in intellect, it seems it should therefore be easily corrected through the proper education. But diligence is a moral character trait, hence connecting stupidity with a characteristic of the will. But there are defects in the will that prevent the elimination of stupidity. There is intellectual pride: a man may be so convinced his opinion is the truth, he remains opaque to any argument based on fact, logic, or metaphysics.

Sloth will make a man too lazy to determine the facts, think logically, or study doctrines. Fear will keep a man from questioning the false beliefs of his family, or social group, or government. So what is the ultimate source of these defects of will?

Death is the other surd in life. Its dominating power occludes any vision of anything beyond our human state in life. It obscures our Being. To stave off death, we feel compelled to satisfy all of life's desires, leaving no room to become conscious of the source of our Being beyond this life.

We become passive toward impulses, instincts, passions, and inclinations, so that anything that might interfere with their satisfaction is rejected. That typically means having to choose between an increase in Life and an increase in Being.

However, we may find ourselves conscious of the source of Being, which is outside and beyond our human life. Being, then, becomes the active force over life's desires, action becomes effortless, much like the lilies in the field.

Why, then, do some receive this grace and others do not? This is a question involving the destiny of every person and is another topic.


\flrightit{Posted on 2011-04-14 by Cologero}

